%%%PAGE 250 rot=0 legibility=good summary=竖直圆轨道脱轨条件、绳碰钉完整圆周、带电粒子是否计重力、圆锥摆与圆锥筒、斜面平抛、斜抛运动公式
%%%BLOCK id=250-01 ch=04 sec=圆周运动·竖直面 kind=problem alt=06 cont=none y=0.00-0.17
%FIG p250_a 0.05 0.0 0.15 0.052 soft
\notefig[0.25]{figures/p250_a.png}
\[
mg(\sin37)^\circ=\frac{mV'^2}{R}
\]
\[
V'^2=\frac35gR
\]
\[
-mg\,\frac85R=\frac12m\left(\frac35gR-V_0^2\right)
\]
\[
V_0^2=\frac{19}{5}gR\qquad V_0=\sqrt{\frac{19}{5}gR}
\]
%%%END
%%%BLOCK id=250-02 ch=04 sec=圆周运动·竖直面 kind=knowledge alt=none cont=none y=0.01-0.12
不脱轨\quad $V_0\le\sqrt{2gR}$\quad $V\ge\sqrt{5gR}$ \\
脱轨\ \unclear{$\Rightarrow$}\ $V_0\in\left(\sqrt{2gR},\sqrt{5gR}\right)$
%%%END
%%%BLOCK id=250-03 ch=04 sec=圆周运动·竖直面 kind=problem alt=06 cont=none y=0.00-0.16
%FIG p250_b 0.57 0.0 0.725 0.12
\notefig[0.3]{figures/p250_b.png}
$v=\sqrt{2gl}$

使其绕 $P$ \unclear{的}完整\unclear{圆周}运动
\[
\sqrt{2gl}\ge\sqrt{5g(l-d)}
\]
%%%END
%%%BLOCK id=250-04 ch=09 sec=带电粒子·是否计重力 kind=note alt=11 cont=none y=0.19-0.25
电子质子原子 % 原稿小字，写在下一行上方 \\
基本粒子、原子核。不计重力\quad 其它计重力
%%%END
%%%BLOCK id=250-05 ch=04 sec=圆周运动·水平面 kind=knowledge alt=none cont=none y=0.28-0.47
\textbf{圆锥摆}
%FIG p250_c 0.065 0.325 0.225 0.468
\notefig[0.25]{figures/p250_c.png}
\[
F_n=mg\tan\theta=\frac{mv^2}{r}=m\omega^2r
\]
\[
r=L\sin\theta\qquad v=\sqrt{gL\sin\theta\tan\theta}\qquad \omega=\sqrt{\frac{g}{L\cos\theta}}=\sqrt{\frac{g}{h}}
\]
若 $\theta\uparrow$ 则 $\omega\uparrow$ $v\uparrow$ $F_n\uparrow$ $T\uparrow$
%%%END
%%%BLOCK id=250-06 ch=04 sec=圆周运动·水平面 kind=knowledge alt=none cont=none y=0.50-0.65
\textbf{圆锥筒}
%FIG p250_d 0.08 0.52 0.30 0.655
\notefig[0.25]{figures/p250_d.png}
\[
\frac{mg}{\tan\theta}=\frac{mv^2}{r}=m\omega^2r
\]
\[
V=\sqrt{\frac{gr}{\tan\theta}}\qquad \omega=\sqrt{\frac{g}{r\tan\theta}}
\]
小球位置越高 $r\uparrow$ $\omega\downarrow$ $v\uparrow$
%%%END
%%%BLOCK id=250-07 ch=04 sec=平抛运动·斜面 kind=knowledge alt=none cont=none y=0.66-0.88
\textbf{平抛斜面}
%FIG p250_e 0.08 0.692 0.285 0.79
\notefig[0.25]{figures/p250_e.png}
位偏角 $\theta$\quad $t=\dfrac{2V_0\tan\theta}{g}$

同角落在斜面 只与 $\theta$ 有关

落回斜面上水平位移与初速度平方成正比

最高点\quad $\dfrac{V_y}{V_x}=\tan\theta$\quad 该刻为全过程中间时刻\quad $t=\dfrac{V_0\tan\theta}{g}$

该点到斜面的最远距离\ $d=\dfrac{V_0^2\sin\theta}{2g\cos\theta}$ % CHECK: 按推导分子似应为 V_0^2 sin^2θ，原稿写 sinθ
%%%END
%%%BLOCK id=250-08 ch=04 sec=斜抛运动 kind=formula alt=none cont=none y=0.88-1.00
\textbf{斜抛运动}
%FIG p250_f 0.05 0.897 0.262 0.97
\notefig[0.25]{figures/p250_f.png}
速度：$V_x=V_{0x}=V_0\cos\theta$\qquad $V_y=V_0\sin\theta=gt$。 % CHECK: 第二个等号原稿为“=”，按物理似应为 V_y=V_0\sin\theta-gt

斜抛运动具有轨迹、速度、运动时间对称，加速度

位移\ $x=V_0\cos\theta\,t$\qquad $y=V_0\sin\theta\,t-\dfrac12gt^2$

射高\ $H=\dfrac{V_{0y}^2}{2g}=\dfrac{(V_0\sin\theta)^2}{2g}$\qquad $t=\dfrac{2V_0\sin\theta}{g}$

水平射程\ $x=V_{0x}t=V_0\cos\theta\,\dfrac{\unclear{2}V_0\sin\theta}{g}$ % 原稿分子前为一小点，按推导应为 2
\[
=\frac{V_0^2\sin2\theta}{g}
\]
%%%END
%%%PAGEEND 250
